% =====================================================================================
% WHITE PAPER introduction (user 2026-11-12): for experimentalists and application engineers. Device and processing
% context first; then what the white paper shows (tilts vs experiment, idealised substrate processing, effect on the band
% gap and on thermal degradation, insights). No model, functional or DFT details. Full paper: overleaf repo.
% =====================================================================================

Halide perovskite solar cells have reached certified single-junction efficiencies above
\SI{26}{\percent}.\cite{ref_efficiency} In a device the absorber is a polycrystalline film, less than a micrometre
thick, sandwiched between an electron- and a hole-transport layer on a glass substrate coated with a transparent
conducting oxide. The tin iodides \CsSnI{}, \MASnI{} and \FASnI{} (MA: methylammonium, FA: formamidinium) are the main
lead-free absorbers, with band gaps of 1.2 to \SI{1.4}{eV}.\cite{pitaro2022tin, stoumpos2013semiconducting} For these
films, the quantities that matter in a device are the band gap at the operating temperature, the crystal phase the film
is in, and how long it stays there. All three depend on the same structural feature, the cooperative tilt of the
SnI$_6$ octahedra, which diffraction detects through superlattice reflections and which changes the gap through the
overlap of the tin and iodine orbitals.\cite{huang2013electronic}

The tilts change close to device temperatures. \CsSnI{} passes from its tetragonal $\beta$ to its orthorhombic $\gamma$
phase at \SI{362}{K}, a few kelvin above the \SI{85}{\celsius} of damp-heat qualification, and at room temperature its
black perovskite phase is metastable against a denser yellow polymorph.\cite{Chung2012JACS} \MASnI{} orders its tilt
near \SI{275}{K},\cite{Yamada2011BCSJ, Takahashi2011Dalton, Parrott2016JPCL, stoumpos2013semiconducting} \FASnI{}
between 225 and \SI{250}{K}.\cite{schueller2018crystal} The compounds also differ in how successive layers of octahedra
tilt: in the same sense in \CsSnI{} and \FASnI{}, in alternating senses in
\MASnI{}.\cite{Chung2012JACS, Kahmann2020ACSEL, stoumpos2013semiconducting}

Manufacturing adds a further variable. The film is coated from solution, crystallised by a thermal anneal and cooled
to room temperature while bonded to its substrate.\cite{StrainReview, yang2022strain} The perovskite contracts several
times more than glass or indium tin oxide on cooling,\cite{Steele2019Science} so the substrate leaves it under an
in-plane tension of tens of MPa, up to about \SI{150}{MPa} after an anneal at \SI{160}{\celsius}, and strains of about
\SI{0.5}{\percent}.\cite{Rolston2018AEM, Zhao2017SciAdv, ref_processing_temp} Tensile films degrade faster than
compressed ones, and the activation energy of their ion migration falls with tension.\cite{Zhao2017SciAdv,
ref_processing_temp} A film that crosses a phase transition while clamped cannot take up the spontaneous strain of the
transition, so the process may also change the structure in which the film ends. Measurements see these effects
together with grain boundaries, interfaces, oxidation and moisture,\cite{ref_sn_instability} and cannot easily isolate
the part that is due to the process.

In this white paper we separate these effects in the three tin iodides with molecular dynamics driven by UrsaMLFF, our
machine-learned force field. We first show that the simulations follow the measured octahedral tilts of \CsSnI{},
\MASnI{} and \FASnI{} and their changes with temperature, including the opposite stacking of the MA and FA compounds.
We then model the coating of a film on a substrate in an idealised form, an anneal followed by cooling with the
in-plane dimensions of the film constrained, and follow its structure to the operating temperature. With this model we
determine how the process shifts the band gap of the film and how the residual strain affects the intrinsic routes of
thermal degradation, iodide migration and, in \CsSnI{}, conversion to the yellow phase. We close with the insights
these results offer for the choice of cation, anneal and substrate.
